\section{Ordenação estável}
A ordenação é considerada estável quando preserva a ordem relativa dos elementos com chaves iguais na sequência de entrada. Em outras palavras, se dois elementos possuem a mesma chave e um deles aparece antes do outro na sequência original, um algoritmo de ordenação estável garantirá que essa ordem relativa seja mantida na sequência ordenada.

A estabilidade em algoritmos de ordenação é especialmente importante em cenários onde os elementos possuem atributos adicionais além da chave utilizada para a ordenação. Por exemplo, em um sistema de gerenciamento de registros, onde cada registro é identificado por um nome e uma data, a estabilidade permite ordenar os registros primeiro por nome e, em seguida, por data, sem perder a ordem relativa estabelecida na primeira ordenação.

A tabela \ref{tab:stable-sort-algorithms} apresenta os algoritmos que serão implementados neste trabalho e indica quais deles são estáveis:

\begin{table}[!h]
    \centering
    \Caption{\label{tab:stable-sort-algorithms}Da estabilidade dos algoritmos apresentados}
    \begin{tabular}{ |p{3cm}|p{3cm}|  }
    \hline
    Estável        & Não estável \\
    \hline
    Bolha          & Seleção     \\
    Bolha Com Flag & Shell Sort  \\
    Coquetel       & Quick Sort  \\
    Inserção       & Heap Sort   \\
    Merge Sort     & Balde       \\
    Contagem       &             \\
    Radix Sort     &             \\
    \hline
    \end{tabular}
\end{table}
